\begin{resumo}
\SingleSpacing\setlength{\parindent}{0pt}
Este trabalho teve como objetivo identificar, mapear e caracterizar depósitos e complexos de transporte de massa no talude norte da Bacia de Santos, com ênfase na análise morfométrica dos eventos. A análise também abordou a distribuição espacial, a organização estratigráfica, a configuração interna, o confinamento e as relações com o relevo do Canal de Santos e a tectônica salífera. A pesquisa utilizou o volume sísmico tridimensional BS-500, com interpretação em amplitude e apoio do atributo cosseno da fase instantânea. Foram mapeadas as superfícies de topo e base dos depósitos, caracterizadas feições associadas à tectônica salífera e identificadas depressões interpretadas como \textit{pockmarks}. As medidas morfométricas foram convertidas do domínio do tempo para profundidade mediante velocidades representativas dos eventos, estimadas a partir de dados de poços e ponderadas pelo inverso da distância, com quantificação das componentes de incerteza consideradas. Complementarmente, aplicou-se a técnica da razão espectral horizontal/vertical a registros de um sismógrafo de fundo oceânico. Foram individualizados dez eventos, distribuídos em três depósitos individuais e dois complexos, estes constituídos por três e quatro eventos, respectivamente. As áreas mapeadas variaram de 11,72 a 1.585,48 km², e os volumes estimados, de 0,444 a 131,6 km³. O relevo associado ao Canal de Santos e à tectônica salífera condicionou a distribuição, o confinamento e os espessamentos dos corpos, enquanto falhas registraram deformação posterior à deposição em parte dos eventos. Foram reconhecidas seis configurações sísmicas internas, com ampla ocorrência de padrões caóticos e coexistência de setores com diferentes graus de organização. As relações estratigráficas envolvendo possíveis \textit{paleopockmarks} e \textit{pockmarks} superficiais indicaram diferentes momentos de formação em relação aos depósitos. A estimativa de profundidade obtida pela razão espectral horizontal/vertical mostrou compatibilidade local com o contato entre os eventos D e C do complexo 2, interpretado como uma interface candidata à resposta de ressonância. Em conjunto, o inventário e a caracterização morfométrica e estratigráfica dos dez eventos permitiram comparar depósitos de diferentes escalas e reconhecer a influência do relevo associado ao Canal de Santos e à tectônica salífera em sua distribuição, seu confinamento e sua acumulação.

\vspace{\onelineskip}
\noindent\textbf{Palavras-chave}: depósitos de transporte de massa; complexos de transporte de massa; sísmica 3D; morfometria; tectônica salífera; Bacia de Santos.
\end{resumo}
